\scanpage{14}
\mspage{2}% typed -2-
literature. The proof of the equivalence of homotopy categories is not in the literature, indeed there is just this M.Phil. thesis of F. Hintze at Warwick which gives full details. (Part of the proof is also given by Federer in notes on algebraic topology from Brown University, 1962.) Cubical sets have a cartesian product, which does not have the correct realisation. This causes all sorts of difficulties. There is a tensor product which has the correct realisation, but having two products certainly complicates the issues. To obtain a Dold-Kan theorem, one has to introduce extra structure on cubical sets, namely the connections considered by Brown-Higgins. Realisations when you have this extra structure have not yet been worked out. More general notions of cubical groups, and related constructions, which have been well worked over by Kan and others in the simplicial case, have not yet been worked out for cubical groups.

$\ell$-3\qquad and \uline{a} corresponding.

p4 $\ell$8\qquad I would prefer here to talk about the Dold-Kan theorem, since Puppe only came in at a later stage, after the theorem had been proved independently by Dold and Kan. In fact Dold's original version was rather computational, and it was Kan who pointed out that the functor from chain complexes to simplicial abelian groups is given in a dimension $n$ by $\operatorname{Hom}(C_{N}\Delta^{n},C)$ for a chain complex $C$ . Also, a description of the theorem for a general additive category was done by Epstein.

It is curious that the way the Loday approach to homotopy type has gone is via $n$-simplicial sets although the proof of the modelling of all homotopy types is as yet incomplete. Similarly, the general van Kampen theory of Brown-Loday uses $n$-fold groupoids, more particularly $\mathrm{Cat}^{n}$-groups . I wonder whether the main problem is that to define gadgets like these for other situations, such as topoi, it may be necessary to have unit interval constructions, and this is where the various stack structures intervene, before obtaining the $\mathrm{Cat}^{n}$-groups.

$\ell$-7\qquad Replace eluding by evading.

$\ell$-2\qquad developed

p5 $\ell$8,9\quad Maybe it would be better as ``whereas it seems to me we have now come to a definite starting point''.

$\ell$11\qquad manifold

p7\qquad This construction of the nerve and its relation to homotopy types has been studied in detail by Fritsch and Latch.

p8\qquad The middle paragraph of this page is interesting. Certainly finding categories which model geometric situations, and then by the nerve obtaining a topological
% archive sheet number: 13
